\section{de Rham Cohomology: Bridging Geometry, Topology, and Physics}
\label{sec:de_rham}

We now return to the reverse question. In ordinary vector calculus, every
gradient field is curl-free. Does every curl-free field come from a globally
defined scalar potential?

Differential forms let us ask the same question on a manifold. Given a
differential $1$-form $\omega$, can we find a smooth function $f$, defined on
the whole manifold, such that
\begin{equation}
    \omega = df?
\end{equation}
Since $d^2 = 0$, such a form must satisfy $d\omega = 0$. We call forms
satisfying this condition \textbf{closed}, and forms that can be written as
$df$ \textbf{exact}. Every exact form is closed, but the converse can fail.

The first \textbf{de Rham cohomology} space records this failure. More
generally, de Rham cohomology asks, at each degree, which closed forms can be
obtained by differentiating a form of one lower degree.

\subsection{Cohomology Groups}
\begin{definition}[de Rham Cohomology Groups]
    For each $k \ge 0$, the \textbf{$k$-th de Rham cohomology group} of $M$ is
    the quotient vector space \cite[p.~441]{lee2013smooth}\footnote{%
        We take $\Omega^{-1}(M)=\{0\}$, so the denominator in degree zero
        is $\{0\}$.}
    \begin{equation}
        H^k_{\mathrm{dR}}(M)
        = \frac{\{\text{closed }k\text{-forms on }M\}}
               {\{\text{exact }k\text{-forms on }M\}}
        = \frac{\ker\!\left(d\colon\Omega^k(M)\to\Omega^{k+1}(M)\right)}
               {\operatorname{im}\!\left(d\colon\Omega^{k-1}(M)\to\Omega^k(M)\right)}.
    \end{equation}
\end{definition}

Since $d$ is linear and $d^2=0$, exact $k$-forms form a linear subspace of
closed $k$-forms. We can therefore define their quotient.

% OPTIONAL BEGIN: Cohomology of Spheres
%\subsection{Cohomology of Spheres}
%While we do not have time in this talk to perform the full calculation, the de
%Rham cohomology groups of the $n$-dimensional sphere $S^n$, for $n \ge 1$,
%are known \cite[Theorem~17.21, p.~450]{lee2013smooth} and illustrate the
%topological meaning of the invariants:
%\begin{equation}
%    H^k_{\mathrm{dR}}(S^n) =
%    \begin{cases}
%        \mathbb{R}, & k = 0 \text{ or } k = n, \\[4pt]
%        0, & 0 < k < n.
%    \end{cases}
%\end{equation}
%These results illustrate how cohomology records different information at
%different degrees. More generally, if two smooth manifolds have different de
%Rham cohomology spaces in some degree, they cannot be diffeomorphic
%\cite[Corollary~17.4, p.~442]{lee2013smooth}.
% OPTIONAL END: Cohomology of Spheres

\subsection{The de Rham Theorem: A Bridge Between Geometry and Topology}
The significance of de Rham cohomology is crystallized by a landmark result in
differential topology (see \cite[Proposition~18.13, pp.~482--483;
Theorem~18.14, p.~484]{lee2013smooth}):

\begin{theorem}[de Rham]
    \label{thm:de_rham}
    For every smooth manifold $M$ and nonnegative integer $k$, there is a
    natural isomorphism of real vector spaces:
    \begin{equation}
        H^k_{\mathrm{dR}}(M) \cong H^k(M; \mathbb{R}),
    \end{equation}
    where $H^k(M; \mathbb{R})$ denotes singular cohomology with real
    coefficients, defined from the topology of $M$ alone.
\end{theorem}

Intuitively, topologists study holes using geometric loops and membranes,
and analogous objects in higher dimensions. Differential forms give us
another way to study them.

The de Rham theorem says that de Rham cohomology recovers exactly the same
information about holes as topological cohomology with real coefficients.
Thus, we can study this topological information through smooth differential
forms and the exterior derivative. The result depends on the topology of the
manifold, even though its definition uses calculus.

\subsection{Physical Significance: Conservative Vector Fields and Topology}
Beyond pure topology, de Rham cohomology plays a decisive role in physics,
explaining theorems that frequently appear mysterious in standard
multivariable calculus courses.

Recall the classical conservative vector field theorem: \emph{on a simply
connected domain $X \subseteq \mathbb{R}^3$, a smooth vector field $F$
satisfies $\operatorname{curl} F = 0$ if and only if $F = \nabla f$ for some
scalar potential $f \colon X \to \mathbb{R}$}
\cite[Corollary~16.27, p.~421]{lee2013smooth}.

\begin{remark}
    If you are unfamiliar with the formal topological definition of ``simply
    connected'', you may simply think of a general well-behaved region in
    $\mathbb{R}^3$ without holes or missing axes (for example, all of
    $\mathbb{R}^3$ or a solid ball).
\end{remark}

Under the dictionary between vector calculus and differential forms:
\begin{enumerate}
    \item $\operatorname{curl} F = 0$ translates to $d\omega_F = 0$, meaning
        that the associated $1$-form $\omega_F$ is \textbf{closed}.
    \item $F = \nabla f$ translates to $\omega_F = df$, meaning that the
        $1$-form is \textbf{exact}.
\end{enumerate}

Thus, the physics theorem that every curl-free field is conservative on $X$ is
precisely the mathematical statement that every closed $1$-form is exact,
which means that the first de Rham cohomology group vanishes:
\begin{equation}
    H^1_{\mathrm{dR}}(X) = 0.
\end{equation}
For a simply connected domain $X$, topology gives $H^1(X; \mathbb{R}) = 0$
\cite[p.~472; Problem~18-2, p.~487]{lee2013smooth}.
The de Rham theorem (Theorem~\ref{thm:de_rham}) then gives
$H^1_{\mathrm{dR}}(X) = 0$, so every closed $1$-form is exact. In
vector-calculus terms, every curl-free field on $X$ has a globally defined
scalar potential. Thus, the question we began with---whether given rates of
change come from one global function---can be answered through topology.
